\chapter{The Tower Turns Around}

\section{The backward descent}

The first three chapters climbed. Each tier was \emph{introduced} atop the one below it --- a class, an inductive, a
constructor, a relation, posited by cases --- and each introduction was a little speculative, which the source
recorded in its own way: the private hedge each type carries rose steadily with the reach, quiet at the bottom and
loudest at the top, where the residue became a sample straining to be a number. This chapter is where the tower turns
around. It stops introducing and starts \emph{eliminating} --- consuming what was built rather than positing more ---
and the first tell is the hedge itself, which does not keep climbing. On the way down it collapses.

Every type in this construction has two halves, and naming them is the cleanest way to see the turn. Following
Gentzen, the rules that \emph{build} an inhabitant of a type are its introduction rules, and the rules that \emph{use}
one are its elimination rules. The first three chapters were all introduction: every tier added something --- a
clause, a constructor, a case --- to the tier beneath it. This section reads the first elimination steps, and the
striking thing is that they recur, in the same shape, at every tier of the tower at once.

The first is the strict order. At no tier is ``strictly less than'' posited on its own; everywhere it is
\emph{derived} from the order already built. The definition is one line, and it is the same line at every rung: strict
precedence is the built order holding one way and failing the other --- \verb|le a b| holding, and \verb|le b a|
failing. It introduces nothing. It takes the \verb|le| the chapter before built and uses it twice, once forward and
once negated, and adds not a single new case. Where each earlier tier had added a clause to the one below, this adds
none; it consumes what is already there. That is the first genuinely backward step in the whole construction --- the
derivation, which had only ever wound upward, beginning to unwind. And it is not a local trick: the very same one-line
elimination stands at the naturals, at the rationals, at the sequences, at the limits, and at the sample, and again at
the trial one episode further on. Six tiers, one elimination, positing nothing.

The second elimination step is quieter and goes further. Having derived its strict order, each tier then \emph{hands
the whole relation away}. A single instance declaration gives the hand-built order to the ambient machinery --- the
\(\leq\) and \(<\) that everything in the surrounding language already speaks --- by pointing the ambient symbol at the
tier's own definition. After that declaration, writing \(a \leq b\) on that tier no longer invokes anything special;
it \emph{means} the relation the tier built by hand. This is elimination in the strict sense: the introduced object is
discharged into the surrounding system and stops being a special case. The tower turns around not by being admired but
by becoming \emph{usable} --- its private order absorbed into the common one, so that the rest of the construction can
speak to it without knowing it was ever bespoke.

And here the hedge tells its second story. On the climb, each type's private doubt rose with how far the reading
reached --- the order rules quiet, the residue louder, the sample loudest of all. On the descent it crashes. The
instance declarations that discharge the order into the ambient one carry the lowest hedges anywhere in the
construction, a near-certainty the source reserves for nothing else. Read that gradient as the chapter's thesis. The
source is \emph{unsure} about what it introduced and \emph{sure} about what it eliminates --- and the asymmetry is
exactly right, because building the tower was a wager and using it is not. Nor is the reversal a single grand arc from
the bottom of the book to here. It recurs, in miniature, at every tier's own turn: the trial's type carries the same
high hedge the sample did --- the wager restated --- its order rule a middling one, and its instance discharge the
same near-certain floor. Introduce, then eliminate; hedge, then be sure; at every rung the same small turn.

In the two verbs, the descent reads cleanly. Introduction was funge: each tier bagged the one below it and added a
clause, gathering by likeness into a larger whole. This section is the first \emph{tange} on the way down. The strict
order \emph{selects} the usable part out of the built one --- forward-and-not-backward, a characteristic picked out of
the relation already in hand. The instance \emph{selects} the bespoke order into the ambient one, choosing the built
relation to stand as the common symbol's meaning. Neither bags anything new. Each is a selection from the bag the
climb already filled --- the tower turning from gathering to picking.

One forward sentence, because it is the sentence the whole book will need at its end. The honest reading, when it
finally comes, will be an \emph{elimination} and not an introduction. The jar the last chapter holds open is a bracket
you \emph{use} --- something to eliminate a claim against, to test a number by --- and never a magnitude you
\emph{posit}, introduced by fiat and defended after. This section plants the reason: the descent, elimination, is the
low-hedge, sound half of every type; the climb, introduction, was the wager. That the tower ``turns around,'' that the
way down is the certain way, that building was a gamble and using is not --- those are the reading, the gloss laid over
a plain fact. The plain fact is smaller and firmer: one line deriving strict order from order, an instance discharging
it into the ambient language, and a hedge that falls exactly as the construction stops guessing and starts spending
what it has.

\section{The seam}

The last section turned the tower around and watched it become usable. This section reads the place where one level
hands to the next --- the \emph{join} between floors --- and finds, at the join, a single thing being carried across:
the residue. The tiers were never a stack of independent floors. Each level is built \emph{over} the one below it, and
what passes upward through every join is the un-eliminable term of the third chapter. The source, climbing, said as
much in an aside --- it was, it admitted, ``just planting seeds.'' This section is where the seeds come up.

In this construction the join is not a new device. It is the instance-implicit constraint itself, read as a seam.
Recall that no tier stands alone: each is defined with the tiers beneath it named as standing conditions --- given a
distinguishable, an admissible, a countable, an encoded, \emph{then} this. That bracketed demand, the place the
construction says ``given all of that, now this,'' \emph{is} the seam between floors. Turning the tower around makes
the seams visible not as demands you climb toward but as joins you descend through.

And here a promise from the third chapter comes due. The residue was called there the top of the tower --- but
carefully, the top of \emph{that} tower, the one the first three chapters climbed. This is where that ``top'' turns out
to be a seam and not a ceiling. A later structure, further up the construction, takes the residue itself as one of its
standing conditions and keeps building --- through more tiers above it, eleven in all, up to the comparable at the
summit, with the physical tier (where the slip's world will enter, in the next chapter) near the top. The residue is a
floor with more floors above it. What the third chapter called a top, this section shows was a landing.

Look at what the higher structure actually binds, and the seam's payload is explicit. Among its standing conditions is
the residue, named as an instance the way the distinguishable and the encoded are named --- the un-eliminable term
carried bodily into a level built over it. And the source, at exactly that binder, muses about what the residue will
turn into on the far side: probably, it guesses, a norm --- a Sobolev norm --- or a cross-product. Those guesses are
not this volume's to cash; they are the reading, a suspicion about the residue's geometric future that the physical
gauge will settle and this one will not. What is firm here is smaller and exact: the residue is the instance the next
level is constructed over. ``The residue at the join'' is not a metaphor. It is the literal payload of the seam --- the
same leftover, passed upward, floor after floor.

The seam is also where a new direction becomes possible, and this is the quiet event the chapter has been walking
toward. To extend a number system by a square root of a negative --- to talk about the imaginary unit at all --- you
need a \emph{direction}, an axis to turn in. And the machine notices, at the seam, that it already has one. The
residue, as the third chapter established, is a direction and a magnitude together. So the direction the imaginary
needs is not conjured; it is the residue's own. The source makes the point in as many words: the complex limit
requires a direction, of which --- it says --- we have a residue already, so it is perfectly cromulent to talk about
\(i\). Read that precisely. This is not yet the rotation; the quarter-turn, the \(-i\) that carries one axis into the
other, is a later chapter's business. This is only the crossing: the first point at which the imaginary becomes
\emph{sayable}, licensed not by fiat but by a direction the construction already built. (The source, in the same
breath, reaches further --- toward a half-integer spin the imaginary will carry --- but that reach is the next
volume's; this one keeps to the axis becoming nameable.)

In the two verbs the seam is a hinge. The lower level was \emph{funged} --- bagged, by likeness, into a tier. At the
join, the next level does not bag again; it \emph{tanges} --- selects the residue out of the tier below and builds over
exactly that. Bag the level, select its residue, seed the next: each seam is a funge-then-tange hinge, and the tower is
a chain of them. The climb was gathering; the joins are where gathering hands to selecting.

One forward sentence. The residue carried across this seam is the same residue the whole instrument will, at its far
end, ask for a number --- the leftover the coupling reads. What this section adds is that the leftover does not sit in
one place; it crosses every join, upward, and arrives at the top still itself. That it will become, on the far side, a
field's geometry, a commutator, a Dirac equation --- those are the readings the source gestures at in its own margins
and hands forward to the physical volume; no such object is built here. This volume builds the seam, carries the
residue through it, and lets the imaginary become sayable on the strength of a direction it already owns. The geometry
the seam is pregnant with is the next book's to deliver.

\section{What came back changed}

The tower climbed, in the first three chapters; it turned around, in the first section of this one; it crossed its
seams, in the last. This section asks the closing question of the chapter --- what came back? --- and answers that it
is not the same thing that went down. The round-trip, the loop chapter promised, keeps the fact. This section names
what the round-trip does \emph{not} keep, and finds that the difference is exactly the residue. The turn-around does
not return you to where you started. It returns you to where you started \emph{plus} the residue --- and that
difference, it will turn out, is the whole point.

Everything the section needs is already built; nothing new is introduced here. It is a single relation, read a second
way. The residue's relation --- the one the antagonist chapter set down, comparing two settled limits --- has, on its
populated arm, three conditions joined together: the two facts equal, the scale of the one no greater than the other,
the index of the one strictly below the other. Read those two settled limits not as ``this one and that one'' but as
``what went down'' and ``what came back,'' and the three conditions fall cleanly into two piles: what the trip kept,
and what the trip changed.

What it kept is the fact. The first condition holds the two facts equal --- and this is nothing new; it is the
lossless core the loop chapter established, the base case that returns a difference to itself unchanged. Whatever came
back is still about the same thing it was sent down to be about. The identity survives the descent whole. If that were
all, the turn-around would be a closed circle, and there would be nothing to measure.

But it is not all. The other two conditions say what changed, and they change in a fixed direction: the scale rises,
the index strictly deepens. And these two are not arbitrary quantities that happened to drift --- they are exactly the
residue's two components, the magnitude and the direction the antagonist chapter read off the very same relation. So
the thing that came back differs from the thing that went down by \emph{precisely the residue}: the same fact, carried
on an advanced scale and a deepened index. ``What came back changed'' is not a vague remark. It is exact --- the pair
of magnitude and direction advanced, while the fact held. The difference between input and output is residue-shaped,
and it could be nothing else.

Here is why the chapter has spent itself getting to this small difference. The difference the turn-around produces
\emph{is} the residue, and the residue is what the rest of the book measures. Nothing downstream reads the thing that
went down, and nothing reads the thing that came back; everything reads the \emph{difference} between them. The next
chapter will meet this difference as a \emph{slip} --- a loop that fails to close by exactly the residue. A later
chapter will bracket it. A later one still will put the coupling-question to it. Each of those is this same difference,
refined one more turn. And so the deepest lesson of the descent is a discipline about reading: the reading is never the
thing itself. The reading is the residue-shaped difference between what you sent down and what came back.

In the two verbs, the round-trip splits as cleanly as its conditions. The fact is \emph{funged} --- kept, the bag
returned intact, member for member. The residue is the \emph{tange} --- the one thing selected out as different, the
displaced part the trip produced. What comes back is the same bag with a member displaced, and the displacement is the
tange the descent made. The chapter ends by handing that displaced tange forward: it is the measurement's raw
material, the difference every later reading will refine.

One boundary, and one forward sentence, and the chapter is closed. The boundary: to call the output a \emph{normal
form}, and the difference a \emph{reduction}, is the reading --- there is no reduction operation in the code, only a
relation whose diagonal, read as a before-and-after, shows the fact held and the residue advanced. Say ``normal form''
if it helps; the firm thing beneath it is the diagonal of a relation already built. The forward sentence: what the
honest ending hands back, at the very end of the book, is exactly this difference and not the thing itself --- not the
magnitude, but a bracket drawn around the residue-shaped difference the turn-around produced. The tower climbed to
build the residue; it turned around to show the residue is the only difference the machine ever returns; and the rest
of the book is the patient reading of that one difference. The chapter closes on the difference. The next opens on what
happens when a loop tries to close over it and cannot.
